%% GENERATED by tools/from_docs.py from stepss-docs. Do not edit by hand:
%% edit the page named below in stepss-docs and regenerate.
%% source: stepss-docs/src/content/docs/models/discrete-controllers.md

\chapter{Discrete controller models}\label{doc:model-dctl}

Discrete controllers (DCTL) in RAMSES implement \textbf{event-driven logic} rather than continuous differential equations. They fire at specific simulation events (voltage crossing a threshold, a timer expiring, a tap position changing) and execute discrete actions such as tripping a generator, shedding load, or adjusting a transformer tap. They run inside RAMSES' event-driven loop and can call \texttt{disp\_disc} to log switching actions.

The RAMSES data keyword is \texttt{DCTL}:

\begin{Verbatim}
DCTL  model_name  controller_name  field1  field2  ... ;
\end{Verbatim}

\begin{docsnote}{Usage in Dynamic Data Files}

Discrete controllers are defined with the \texttt{DCTL} keyword as standalone records in the dynamic data file. The syntax is:

\begin{Verbatim}
DCTL model_name ctl_name param1 param2 ... ;
\end{Verbatim}

Where \texttt{model\_name} is the controller type, \texttt{ctl\_name} is a unique instance name, and the remaining parameters are model-specific.

Recognised discrete-controller model names (case-sensitive):

\begin{itemize}
\tightlist
\item
  \textbf{Uppercase short names} (no prefix): \texttt{PST}, \texttt{LTC}, \texttt{LTC2}, \texttt{OLTC2}, \texttt{LTCINV}, \texttt{MAIS}, \texttt{UVLS}, \texttt{RT}, \texttt{UVPROT}, \texttt{FRT}, \texttt{VOLT\_VAR}, \texttt{SIM\_MINMAXVOLT}, \texttt{SIM\_MINMAXSPEED}.
\item
  \textbf{Prefixed name}: RAMSES adds the \texttt{dctl\_} prefix automatically, so both \texttt{line\_prot} and \texttt{dctl\_line\_prot} resolve to the line-protection model.
\end{itemize}

All of the above are built into every RAMSES distribution (standalone executable and shared library used by stepss). \texttt{MAIS} is an automatic shunt-reactor switching scheme (\emph{Manoeuvre Automatique d'Inductances Shunt}) that switches susceptance at a monitored bus on under-voltage, voltage-drop and over-voltage thresholds; it does not have a dedicated documentation section below. \texttt{HVDC\_LIM}, \texttt{injprot} and \texttt{LOSPROT}, documented below, have been callable since 3.79; before that they were compiled under no name.

\end{docsnote}

\begin{center}\rule{0.5\linewidth}{0.5pt}\end{center}

\section{Tap Changers}\label{model-dctl:tap-changers}

\begin{docsnote}{TRFO and DCTL LTC}

The \hyperref[ch-network]{TRFO record} combines the transformer model with load tap changer data (tap range, positions, controlled bus, voltage setpoint), but this data is used \textbf{only by the power flow} for ratio adjustment. To control a transformer's ratio during dynamic simulation with RAMSES, a DCTL LTC controller (below) must be associated with the transformer, whether it was defined with TRANSFO or TRFO.

\end{docsnote}

\subsection{\texorpdfstring{LTC (\texttt{dctl\_ltc}): Load Tap Changer (Standard)}{LTC (dctl\_ltc): Load Tap Changer (Standard)}}\label{model-dctl:ltc-dctl_ltc-load-tap-changer-standard}

\subsubsection{Description}\label{model-dctl:description}

Controls the ratio of a transformer to regulate voltage at a \textbf{monitored bus}. It monitors the voltage, compares it against a deadband around the setpoint (initialised to the voltage at \(t=0\)), and initiates a tap change after a delay if the voltage remains outside the deadband.

The first tap change uses a longer delay (\texttt{delay\_first}); subsequent changes use a shorter delay (\texttt{delay\_next}), implementing the standard IEC/IEEE slow-start behaviour.

\subsubsection{Logic Description}\label{model-dctl:logic-description}

Let \(V\) be the monitored bus voltage and \(V_{set}\) the setpoint:

\[
\text{Error} = V_{set} - V
\]

\textbf{Deadband check}: No action if \(|V - V_{set}| \le \delta V\) (half-deadband).

\textbf{Direction logic}: If \texttt{direction\ \textgreater{}\ 0}, increasing the ratio increases the voltage.

\textbf{Timer logic}:

\[
\text{if } |V - V_{set}| > \delta V \text{ and } (t - t_{last}) \ge \tau_{delay}:
\quad \text{tap} \leftarrow \text{tap} \pm \Delta r, \quad t_{last} \leftarrow t
\]

where \(\tau_{delay} = \tau_{first}\) for the first operation, \(\tau_{next}\) thereafter, and the tap is clamped to \([r_{\min}, r_{\max}]\).

\textbf{State variable} \texttt{w(13)}:
- \texttt{0}: voltage within deadband
- \texttt{+1}: \(V < V_{set} - \delta V\) (tap must increase)
- \texttt{\ensuremath{-}1}: \(V > V_{set} + \delta V\) (tap must decrease)

\subsubsection{Parameters}\label{model-dctl:parameters}

{\def\LTcaptype{none} % do not increment counter
\begin{small}\begin{longtable}[]{@{}
  >{\raggedright\arraybackslash}p{(\linewidth - 4\tabcolsep) * \real{0.1892}}
  >{\raggedright\arraybackslash}p{(\linewidth - 4\tabcolsep) * \real{0.4595}}
  >{\raggedright\arraybackslash}p{(\linewidth - 4\tabcolsep) * \real{0.3514}}@{}}
\toprule\noalign{}
\begin{minipage}[b]{\linewidth}\raggedright
Field
\end{minipage} & \begin{minipage}[b]{\linewidth}\raggedright
Working variable
\end{minipage} & \begin{minipage}[b]{\linewidth}\raggedright
Description
\end{minipage} \\
\midrule\noalign{}
\endhead
\bottomrule\noalign{}
\endlastfoot
\texttt{branch\_name} & \texttt{w(1)} & Transformer branch name (must be type \texttt{trfo}) \\
\texttt{bus\_name} & \texttt{w(2)} & Monitored (controlled) bus name \\
\texttt{direction} & \texttt{w(3)} & Sign convention: \textgreater0 means increasing ratio raises voltage \\
\texttt{ratio\_min} & \texttt{w(4)} & Minimum transformer ratio (\textbf{\% of nominal}, e.g.~\texttt{85.} for 0.85 pu) \\
\texttt{ratio\_max} & \texttt{w(5)} & Maximum transformer ratio (\textbf{\% of nominal}, e.g.~\texttt{115.} for 1.15 pu) \\
\texttt{nb\_positions} & \texttt{w(6)} & Number of tap positions (step = (max\ensuremath{-}min)/(nb\ensuremath{-}1)) \\
\texttt{half\_deadband} & \texttt{w(7)} & Voltage half-deadband \(\delta V\) (pu) \\
\texttt{delay\_first} & \texttt{w(8)} & Delay before first tap change (s) \\
\texttt{delay\_next} & \texttt{w(9)} & Delay between subsequent tap changes (s) \\
\end{longtable}\end{small}
}

\begin{docsnote}{Caution}

The ratio limits are given in \textbf{percent}, not per unit: the model divides \texttt{ratio\_min} and \texttt{ratio\_max} by 100 when reading the data. A record with \texttt{0.85\ 1.15} would silently produce limits of 0.0085/0.0115 pu.

\end{docsnote}

\subsubsection{Usage Example}\label{model-dctl:usage-example}

\begin{Verbatim}
DCTL  LTC  LTC_TR1  TR1-HV-MV  BUS_MV  1  85.  115.  33  0.01  30.0  10.0 ;
\end{Verbatim}

\begin{center}\rule{0.5\linewidth}{0.5pt}\end{center}

\subsection{\texorpdfstring{LTC2 (\texttt{dctl\_ltc2}): Load Tap Changer (Variant 2)}{LTC2 (dctl\_ltc2): Load Tap Changer (Variant 2)}}\label{model-dctl:ltc2-dctl_ltc2-load-tap-changer-variant-2}

\subsubsection{Description}\label{model-dctl:description-1}

Identical logic to \texttt{dctl\_ltc} but the \textbf{voltage setpoint is an explicit input parameter} rather than being taken from the initial operating point. This allows prescribing a setpoint different from the initial bus voltage, which is useful when the initial load-flow voltage differs from the desired operating voltage.

\subsubsection{Logic Description}\label{model-dctl:logic-description-1}

Same deadband, direction, and timer logic as \texttt{dctl\_ltc}. The voltage setpoint \(V_{set}\) is read directly from the data field \texttt{Vset} instead of being captured at \(t = 0\).

\[
\text{if } |V - V_{set}| > \delta V \text{ and } (t - t_{last}) \ge \tau_{delay}: \quad \text{tap change}
\]

\subsubsection{Parameters}\label{model-dctl:parameters-1}

{\def\LTcaptype{none} % do not increment counter
\begin{small}\begin{longtable}[]{@{}lll@{}}
\toprule\noalign{}
Field & Working variable & Description \\
\midrule\noalign{}
\endhead
\bottomrule\noalign{}
\endlastfoot
\texttt{branch\_name} & \texttt{w(1)} & Transformer branch (type \texttt{trfo}) \\
\texttt{bus\_name} & \texttt{w(2)} & Monitored bus \\
\texttt{direction} & \texttt{w(3)} & \textgreater0: ratio increase raises voltage \\
\texttt{ratio\_min} & \texttt{w(4)} & Minimum ratio (\% of nominal) \\
\texttt{ratio\_max} & \texttt{w(5)} & Maximum ratio (\% of nominal) \\
\texttt{nb\_positions} & \texttt{w(6)} & Number of tap positions \\
\texttt{half\_deadband} & \texttt{w(7)} & Voltage half-deadband (pu) \\
\texttt{Vset} & \texttt{w(8)} & Voltage setpoint (pu) \\
\texttt{delay\_first} & \texttt{w(9)} & First tap change delay (s) \\
\texttt{delay\_next} & \texttt{w(10)} & Subsequent tap change delay (s) \\
\end{longtable}\end{small}
}

\subsubsection{Usage Example}\label{model-dctl:usage-example-1}

\begin{Verbatim}
DCTL  LTC2  LTC2_TR1  TR1-HV-MV  BUS_MV  -1  88.  120.  33  0.01  1.0  30  8 ;
\end{Verbatim}

\begin{center}\rule{0.5\linewidth}{0.5pt}\end{center}

\subsection{\texorpdfstring{LTCINV (\texttt{dctl\_ltcinv}): Inverse-Time Load Tap Changer}{LTCINV (dctl\_ltcinv): Inverse-Time Load Tap Changer}}\label{model-dctl:ltcinv-dctl_ltcinv-inverse-time-load-tap-changer}

\subsubsection{Description}\label{model-dctl:description-2}

An \textbf{inverse-time} tap changer in which the delay before a tap change is proportional to the integral of the voltage deviation rather than a fixed timer. Large voltage errors trigger faster operation; small but persistent deviations trigger slower operation.

This implements behaviour analogous to inverse-time overcurrent relays, applied to voltage control.

\subsubsection{Logic Description}\label{model-dctl:logic-description-2}

The controller integrates the signed voltage deviation over time:

\[
\Sigma = \int_{t_0}^{t} (V_{set} - V) \, dt'
\]

A tap change fires when the integrated error exceeds a threshold derived from the maximum delay and acceleration parameter:

\[
\text{if } |\Sigma| \ge \frac{\tau_{max} \cdot \delta V}{\text{accel}}: \quad \text{tap change}, \quad \Sigma \leftarrow 0
\]

where \(\tau_{max}\) is the maximum delay and \texttt{accel} is the acceleration factor (smaller value = faster response).

\textbf{State variable} \texttt{w(13)}: same coding as \texttt{dctl\_ltc}.

\subsubsection{Parameters}\label{model-dctl:parameters-2}

{\def\LTcaptype{none} % do not increment counter
\begin{small}\begin{longtable}[]{@{}lll@{}}
\toprule\noalign{}
Field & Working variable & Description \\
\midrule\noalign{}
\endhead
\bottomrule\noalign{}
\endlastfoot
\texttt{branch\_name} & \texttt{w(1)} & Transformer branch \\
\texttt{bus\_name} & \texttt{w(2)} & Monitored bus \\
\texttt{direction} & \texttt{w(3)} & Sign convention \\
\texttt{ratio\_min} & \texttt{w(4)} & Minimum ratio (\% of nominal) \\
\texttt{ratio\_max} & \texttt{w(5)} & Maximum ratio (\% of nominal) \\
\texttt{nb\_positions} & \texttt{w(6)} & Number of tap positions \\
\texttt{half\_deadband} & \texttt{w(7)} & Voltage half-deadband (pu) \\
\texttt{Vset} & \texttt{w(8)} & Voltage setpoint (pu) \\
\texttt{delay\_max} & \texttt{w(9)} & Maximum delay before tap change (s) \\
\texttt{accel} & \texttt{w(10)} & Acceleration factor (s), smaller = faster \\
\end{longtable}\end{small}
}

\subsubsection{Usage Example}\label{model-dctl:usage-example-2}

\begin{Verbatim}
DCTL  LTCINV  LTCINV_TR1  TR1-HV-MV  BUS_MV  1  85.  115.  33  0.01  1.025  60.0  0.1 ;
\end{Verbatim}

\begin{center}\rule{0.5\linewidth}{0.5pt}\end{center}

\subsection{\texorpdfstring{OLTC2 (\texttt{dctl\_oltc2}): On-Load Tap Changer Type 2}{OLTC2 (dctl\_oltc2): On-Load Tap Changer Type 2}}\label{model-dctl:oltc2-dctl_oltc2-on-load-tap-changer-type-2}

\subsubsection{Description}\label{model-dctl:description-3}

An advanced on-load tap changer with \textbf{R/X line drop compensation}, end-stop flags. It includes separate flags to prevent tap movement when the transformer is at its mechanical end-stop, and applies voltage correction based on the current flowing through the transformer using the compensation resistances \texttt{R\_comp} and \texttt{X\_comp}.

The \textbf{Ratio2-factor} enables a secondary correction factor for the transformer ratio independent of the main tap position.

\subsubsection{Logic Description}\label{model-dctl:logic-description-3}

Compensated voltage seen by the controller:

\[
V_{comp} = V_{meas} - R_{comp} \, I_{re} - X_{comp} \, I_{im}
\]

where \(I_{re}, I_{im}\) are the real and imaginary components of the transformer current. The deadband check and timer logic then operate on \(V_{comp}\).

End-stop flags \texttt{w(15)} and \texttt{w(16)} prevent further tap changes when the ratio is already at its limit and the requested change would move it beyond.

\subsubsection{Parameters}\label{model-dctl:parameters-3}

{\def\LTcaptype{none} % do not increment counter
\begin{small}\begin{longtable}[]{@{}lll@{}}
\toprule\noalign{}
Field & Working variable & Description \\
\midrule\noalign{}
\endhead
\bottomrule\noalign{}
\endlastfoot
\texttt{branch\_name} & \texttt{w(1)} & Transformer branch \\
\texttt{bus\_name} & \texttt{w(2)} & Controlled bus \\
\texttt{direction} & \texttt{w(3)} & Sign convention \\
\texttt{ratio\_min} & \texttt{w(4)} & Minimum ratio (pu) \\
\texttt{ratio\_max} & \texttt{w(5)} & Maximum ratio (pu) \\
\texttt{nb\_positions} & \texttt{w(6)} & Number of tap positions \\
\texttt{half\_deadband} & \texttt{w(7)} & Voltage half-deadband (pu) \\
\texttt{Vset} & \texttt{w(8)} & Voltage setpoint (pu) \\
\texttt{delay\_first} & \texttt{w(9)} & First tap change delay (s) \\
\texttt{delay\_next} & \texttt{w(10)} & Subsequent tap change delay (s) \\
\texttt{R\_comp} & \texttt{w(17)} & Line drop compensation resistance (pu) \\
\texttt{X\_comp} & \texttt{w(18)} & Line drop compensation reactance (pu) \\
\texttt{ratio2\_factor} & \texttt{w(20)} & Secondary ratio correction factor (pu) \\
\end{longtable}\end{small}
}

14 fields required (includes from-bus, to-bus, circuit ID for PSS/E format).

\subsubsection{Usage Example}\label{model-dctl:usage-example-3}

\begin{Verbatim}
DCTL  OLTC2  OLTC_TR2  FROM_BUS  TO_BUS  1  1  0.88  1.12  25  0.0075
             1.02  45.0  15.0  0.0  0.0 ;
\end{Verbatim}

\begin{center}\rule{0.5\linewidth}{0.5pt}\end{center}

\section{Protection}\label{model-dctl:protection}

\subsection{\texorpdfstring{UVPROT (\texttt{dctl\_uvprot}): Under-Voltage Protection Relay}{UVPROT (dctl\_uvprot): Under-Voltage Protection Relay}}\label{model-dctl:uvprot-dctl_uvprot-under-voltage-protection-relay}

\subsubsection{Description}\label{model-dctl:description-4}

Trips a \textbf{generator or injector} when the voltage at a monitored bus falls below a threshold for a sustained time delay. This models definite-time under-voltage protection relays.

\subsubsection{Logic Description}\label{model-dctl:logic-description-4}

The relay monitors voltage \(V_{meas}\) at a specified bus:

\[
\text{if } V_{meas} < V_{\min}: \quad \text{start timer}
\]
\[
\text{if timer} \ge T_{delay}: \quad \text{trip generator}, \quad \text{state} \leftarrow -1 \text{ (latched)}
\]

Once tripped (\texttt{w(8)\ =\ -1}), the controller is permanently latched and takes no further action.

\subsubsection{Parameters}\label{model-dctl:parameters-4}

{\def\LTcaptype{none} % do not increment counter
\begin{small}\begin{longtable}[]{@{}lll@{}}
\toprule\noalign{}
Field & Working variable & Description \\
\midrule\noalign{}
\endhead
\bottomrule\noalign{}
\endlastfoot
\texttt{bus\_name} & \texttt{w(1)} & Monitored bus name \\
\texttt{gen\_name} & \texttt{w(2)} & Generator/injector to trip \\
\texttt{Vmin} & \texttt{w(3)} & Under-voltage pickup threshold (pu) \\
\texttt{T\_delay} & \texttt{w(4)} & Time delay before tripping (s) \\
\end{longtable}\end{small}
}

Internal variables: \texttt{w(5)} = measured voltage, \texttt{w(6)} = start time of violation, \texttt{w(7)} = simulation time, \texttt{w(8)} = controller state.

\subsubsection{Usage Example}\label{model-dctl:usage-example-4}

\begin{Verbatim}
DCTL  UVPROT  UV_GEN1  BUS_HV  GEN1  0.85  0.5 ;
\end{Verbatim}

\begin{center}\rule{0.5\linewidth}{0.5pt}\end{center}

\subsection{\texorpdfstring{UVLS (\texttt{dctl\_uvls}): Under-Voltage Load Shedding}{UVLS (dctl\_uvls): Under-Voltage Load Shedding}}\label{model-dctl:uvls-dctl_uvls-under-voltage-load-shedding}

\subsubsection{Description}\label{model-dctl:description-5}

Sheds load in \textbf{multiple steps} in response to sustained voltage depression. The controller integrates the area between the measured voltage and a threshold (energy integral criterion) and triggers load shedding when the area exceeds a limit \(C\). This prevents unnecessary shedding for short transient dips.

After each shedding step, the controller resets the integral and waits a minimum delay before the next step.

\subsubsection{Logic Description}\label{model-dctl:logic-description-5}

Area integral (trapezoidal):

\[
A(t) = \int_{t_{start}}^{t} \max(V_{th} - V_{meas}, 0) \, dt'
\]

Shedding condition:

\[
\text{if } A \ge C \text{ and } (t - t_{last}) \ge T_{\min}: \quad \text{shed load step}
\]

Amount shed per step (bounded):

\[
\Delta P_{shed} = \text{clip}\left(K \cdot A, \; P_{\min,step}, \; P_{\max,step}\right)
\]

Controller terminates when total shedding reaches the fraction \texttt{frac\_load} of initial load or when \texttt{max\_steps} is exhausted.

\subsubsection{Parameters}\label{model-dctl:parameters-5}

{\def\LTcaptype{none} % do not increment counter
\begin{small}\begin{longtable}[]{@{}
  >{\raggedright\arraybackslash}p{(\linewidth - 4\tabcolsep) * \real{0.1892}}
  >{\raggedright\arraybackslash}p{(\linewidth - 4\tabcolsep) * \real{0.4595}}
  >{\raggedright\arraybackslash}p{(\linewidth - 4\tabcolsep) * \real{0.3514}}@{}}
\toprule\noalign{}
\begin{minipage}[b]{\linewidth}\raggedright
Field
\end{minipage} & \begin{minipage}[b]{\linewidth}\raggedright
Working variable
\end{minipage} & \begin{minipage}[b]{\linewidth}\raggedright
Description
\end{minipage} \\
\midrule\noalign{}
\endhead
\bottomrule\noalign{}
\endlastfoot
\texttt{bus\_name} & \texttt{w(1)} & Monitored bus \\
\texttt{load\_name} & \texttt{w(2)} & Controlled load (injector) \\
\texttt{V\_th} & \texttt{w(3)} & Voltage threshold (pu) \\
\texttt{C} & \texttt{w(4)} & Maximum area \(\int(V_{th}-V) \, dt\) before shedding (pu\ensuremath{\cdot}s) \\
\texttt{T\_min} & \texttt{w(5)} & Minimum delay between steps (s) \\
\texttt{T\_max} & \texttt{w(6)} & Maximum delay between steps (s) \\
\texttt{frac\_load} & \texttt{w(7)} & Maximum fraction of load that may be shed \\
\texttt{K} & \texttt{w(8)} & Gain: area \ensuremath{\rightarrow} shed amount \\
\texttt{P\_min\_step} & \texttt{w(9)} & Minimum load shedding per step (MW) \\
\texttt{P\_max\_step} & \texttt{w(10)} & Maximum load shedding per step (MW) \\
\texttt{max\_steps} & \texttt{w(11)} & Maximum number of shedding steps \\
\end{longtable}\end{small}
}

Internal variables: \texttt{w(12)}--\texttt{w(18)} (previous voltage, timers, area accumulator, step counter, total shed).

\subsubsection{Usage Example}\label{model-dctl:usage-example-5}

\begin{Verbatim}
DCTL  UVLS  UVLS1  BUS_MV  LOAD1  0.90  0.5  5.0  20.0  0.3  2.0  5.0  50.0  5 ;
\end{Verbatim}

\begin{center}\rule{0.5\linewidth}{0.5pt}\end{center}

\subsection{\texorpdfstring{\texttt{FRT} (\texttt{dctl\_FRT}): Fault Ride-Through Controller}{FRT (dctl\_FRT): Fault Ride-Through Controller}}\label{model-dctl:frt-dctl_frt-fault-ride-through-controller}

\subsubsection{Description}\label{model-dctl:description-6}

Implements a \textbf{piecewise-linear voltage--time FRT envelope} for an inverter-based generator (IBG). It monitors the bus voltage of the controlled injector and trips the unit if the voltage stays below the FRT envelope longer than allowed by the characteristic.

The FRT envelope is defined by four (voltage, time) breakpoints: \((V_1, t_1), (V_2, t_2), (V_3, t_3), (V_4)\). If the voltage is above the envelope at all times, no action is taken. If it falls below the envelope for more than the corresponding time limit, the unit is tripped.

\subsubsection{Logic Description}\label{model-dctl:logic-description-6}

The FRT curve defines the minimum voltage \(V(t)\) that must be sustained:

{\def\LTcaptype{none} % do not increment counter
\begin{small}\begin{longtable}[]{@{}ll@{}}
\toprule\noalign{}
Segment & Condition \\
\midrule\noalign{}
\endhead
\bottomrule\noalign{}
\endlastfoot
\(V \ge V_1\) & Always ride-through \\
\(V_2 \le V < V_1\) & Must recover within \(t_1\) s \\
\(V_3 \le V < V_2\) & Must recover within \(t_2\) s \\
\(V_4 \le V < V_3\) & Must recover within \(t_3\) s \\
\(V < V_4\) & Immediate trip \\
\end{longtable}\end{small}
}

Two internal timers (\texttt{w(8)}, \texttt{w(9)}) track time spent in violation zones.

\subsubsection{Parameters}\label{model-dctl:parameters-6}

{\def\LTcaptype{none} % do not increment counter
\begin{small}\begin{longtable}[]{@{}
  >{\raggedright\arraybackslash}p{(\linewidth - 4\tabcolsep) * \real{0.1892}}
  >{\raggedright\arraybackslash}p{(\linewidth - 4\tabcolsep) * \real{0.4595}}
  >{\raggedright\arraybackslash}p{(\linewidth - 4\tabcolsep) * \real{0.3514}}@{}}
\toprule\noalign{}
\begin{minipage}[b]{\linewidth}\raggedright
Field
\end{minipage} & \begin{minipage}[b]{\linewidth}\raggedright
Working variable
\end{minipage} & \begin{minipage}[b]{\linewidth}\raggedright
Description
\end{minipage} \\
\midrule\noalign{}
\endhead
\bottomrule\noalign{}
\endlastfoot
\texttt{inj\_name} & \texttt{w(10)} & Controlled injector \\
\texttt{Val1} & \texttt{w(1)} & Voltage threshold 1 (pu), upper FRT boundary \\
\texttt{t1} & \texttt{w(2)} & Time limit for zone 1 (s) \\
\texttt{Val2} & \texttt{w(3)} & Voltage threshold 2 (pu) \\
\texttt{t2} & \texttt{w(4)} & Time limit for zone 2 (s) \\
\texttt{Val3} & \texttt{w(5)} & Voltage threshold 3 (pu) \\
\texttt{t3} & \texttt{w(6)} & Time limit for zone 3 (s) \\
\texttt{Val4} & \texttt{w(7)} & Voltage threshold 4 (pu), below this: trip immediately \\
\end{longtable}\end{small}
}

8 fields required. Field 1 is the injector name, fields 2--8 map to \texttt{Val1,\ t1,\ Val2,\ t2,\ Val3,\ t3,\ Val4}.

\subsubsection{Usage Example}\label{model-dctl:usage-example-6}

\begin{Verbatim}
DCTL  FRT  FRT_WTG1  WT1  0.90  0.15  0.70  0.5  0.20  1.0  0.05 ;
\end{Verbatim}

\begin{center}\rule{0.5\linewidth}{0.5pt}\end{center}

\subsection{dctl\_injprot: Injector Protection}\label{model-dctl:dctl_injprot-injector-protection}

\begin{docsnote}{Note}

Callable since 3.79. Earlier releases compiled this model under no name, so a data file referring to it was rejected.

\end{docsnote}

\subsubsection{Description}\label{model-dctl:description-7}

A \textbf{generic injector protection} relay that monitors any observable variable of an injector (e.g., current magnitude, active power, reactive power) and trips the injector if the variable goes outside prescribed bounds for a sustained time.

The controller identifies the variable to monitor by name from the injector's observable list, making it flexible and applicable to any injector model.

\subsubsection{Logic Description}\label{model-dctl:logic-description-7}

\[
\text{if } x_{obs} < L_{lower} \text{ or } x_{obs} > L_{upper}: \quad \text{start timer}
\]
\[
\text{if timer} \ge T_{trip}: \quad \text{trip injector}, \quad \text{state} \leftarrow -1
\]

The controller latches once tripped and cannot reset. State variable \texttt{w(7)}:
- \texttt{1} = active (monitoring)
- \texttt{0} = idle (within bounds)
- \texttt{\ensuremath{-}1} = already tripped

\subsubsection{Parameters}\label{model-dctl:parameters-7}

{\def\LTcaptype{none} % do not increment counter
\begin{small}\begin{longtable}[]{@{}
  >{\raggedright\arraybackslash}p{(\linewidth - 4\tabcolsep) * \real{0.1892}}
  >{\raggedright\arraybackslash}p{(\linewidth - 4\tabcolsep) * \real{0.4595}}
  >{\raggedright\arraybackslash}p{(\linewidth - 4\tabcolsep) * \real{0.3514}}@{}}
\toprule\noalign{}
\begin{minipage}[b]{\linewidth}\raggedright
Field
\end{minipage} & \begin{minipage}[b]{\linewidth}\raggedright
Working variable
\end{minipage} & \begin{minipage}[b]{\linewidth}\raggedright
Description
\end{minipage} \\
\midrule\noalign{}
\endhead
\bottomrule\noalign{}
\endlastfoot
\texttt{inj\_name} & \texttt{w(1)} & Monitored/controlled injector name \\
\texttt{var\_name} & \texttt{w(2)} & Observable variable name (from injector's observable list) \\
\texttt{lower\_bound} & \texttt{w(3)} & Lower bound for the variable \\
\texttt{upper\_bound} & \texttt{w(4)} & Upper bound for the variable \\
\texttt{T\_trip} & \texttt{w(5)} & Time out-of-bounds before tripping (s) \\
\end{longtable}\end{small}
}

5 fields required.

\subsubsection{Usage Example}\label{model-dctl:usage-example-7}

\begin{Verbatim}
DCTL  injprot  PROT_WTG1  WT1  I  0.0  1.2  0.1 ;
\end{Verbatim}

\begin{center}\rule{0.5\linewidth}{0.5pt}\end{center}

\subsection{dctl\_losprot: Loss-of-Synchronism Protection}\label{model-dctl:dctl_losprot-loss-of-synchronism-protection}

\begin{docsnote}{Note}

Callable since 3.79. Earlier releases compiled this model under no name, so a data file referring to it was rejected.

\end{docsnote}

\subsubsection{Description}\label{model-dctl:description-8}

Trips a \textbf{synchronous generator} when its rotor speed exceeds a maximum threshold for a sustained time, detecting loss of synchronism (pole slipping). The controller monitors the machine's per-unit rotor speed \(\omega\) and triggers disconnection when \(\omega > \omega_{max}\) for longer than the specified time delay.

\subsubsection{Logic Description}\label{model-dctl:logic-description-8}

\[
\text{if } \omega > \omega_{max}: \quad \text{start timer}
\]
\[
\text{if timer} \ge T_{delay}: \quad \text{trip generator}, \quad \text{state} \leftarrow -1
\]

The timer resets if the speed returns within bounds. The controller latches permanently after tripping.

\subsubsection{Parameters}\label{model-dctl:parameters-8}

{\def\LTcaptype{none} % do not increment counter
\begin{small}\begin{longtable}[]{@{}lll@{}}
\toprule\noalign{}
Field & Working variable & Description \\
\midrule\noalign{}
\endhead
\bottomrule\noalign{}
\endlastfoot
\texttt{gen\_name} & \texttt{w(1)} & Controlled synchronous generator name \\
\texttt{omega\_max} & \texttt{w(2)} & Maximum rotor speed threshold (pu) \\
\texttt{T\_delay} & \texttt{w(3)} & Time delay before tripping (s) \\
\end{longtable}\end{small}
}

3 fields required. Internal variables: \texttt{w(4)} = measured speed, \texttt{w(5)} = violation start time, \texttt{w(6)} = current time, \texttt{w(7)} = state.

\subsubsection{Usage Example}\label{model-dctl:usage-example-8}

\begin{Verbatim}
DCTL  LOSPROT  LOS_GEN2  GEN2  1.05  0.2 ;
\end{Verbatim}

\begin{center}\rule{0.5\linewidth}{0.5pt}\end{center}

\subsection{\texorpdfstring{line\_prot (\texttt{dctl\_line\_prot}): Line Overcurrent Protection}{line\_prot (dctl\_line\_prot): Line Overcurrent Protection}}\label{model-dctl:line_prot-dctl_line_prot-line-overcurrent-protection}

The data-file model name is \texttt{line\_prot} (or \texttt{dctl\_line\_prot}). RAMSES adds the \texttt{dctl\_} prefix automatically.

\subsubsection{Description}\label{model-dctl:description-9}

Monitors the \textbf{apparent power flow on a set of lines} and trips them if they exceed their thermal rating (with a configurable oversize factor). It is designed to protect multiple lines simultaneously and disconnects any line whose flow exceeds \texttt{oversize\_factor\ \ensuremath{\times}\ S\_max}.

The controller uses RAMSES' internal \texttt{smax\_bra} array (maximum apparent power ratings from the network data) and reads real-time power flows via \texttt{pqbra}.

\subsubsection{Logic Description}\label{model-dctl:logic-description-9}

For each monitored branch \(k\):

\[
\text{if } |S_k| > \text{oversize\_factor} \times S_{max,k}: \quad \text{disconnect branch } k
\]

The action is immediate (no time delay); the oversize factor allows temporary overloads.

\subsubsection{Parameters}\label{model-dctl:parameters-9}

{\def\LTcaptype{none} % do not increment counter
\begin{small}\begin{longtable}[]{@{}
  >{\raggedright\arraybackslash}p{(\linewidth - 4\tabcolsep) * \real{0.1892}}
  >{\raggedright\arraybackslash}p{(\linewidth - 4\tabcolsep) * \real{0.4595}}
  >{\raggedright\arraybackslash}p{(\linewidth - 4\tabcolsep) * \real{0.3514}}@{}}
\toprule\noalign{}
\begin{minipage}[b]{\linewidth}\raggedright
Field
\end{minipage} & \begin{minipage}[b]{\linewidth}\raggedright
Working variable
\end{minipage} & \begin{minipage}[b]{\linewidth}\raggedright
Description
\end{minipage} \\
\midrule\noalign{}
\endhead
\bottomrule\noalign{}
\endlastfoot
\texttt{oversize\_factor} & \texttt{w(1)} & Allowed overload factor (e.g., 1.05 = 5\% overload permitted) \\
\texttt{nb\_lines} & \texttt{w(2)} & Number of protected lines \\
\texttt{line\_1}, \ldots{} & \texttt{w(3-)} & Names of protected branches (space-separated) \\
\end{longtable}\end{small}
}

\subsubsection{Usage Example}\label{model-dctl:usage-example-9}

\begin{Verbatim}
DCTL  line_prot  LINEPROT1  1.05  2  L1-L2  L3-L4 ;
\end{Verbatim}

\begin{center}\rule{0.5\linewidth}{0.5pt}\end{center}

\begin{center}\rule{0.5\linewidth}{0.5pt}\end{center}

\subsection{\texorpdfstring{\texttt{PST} (\texttt{dctl\_pst}): Phase-Shifting Transformer Control}{PST (dctl\_pst): Phase-Shifting Transformer Control}}\label{model-dctl:pst-dctl_pst-phase-shifting-transformer-control}

\subsubsection{Description}\label{model-dctl:description-10}

Controls the phase angle (and thus the \textbf{active power flow}) of a phase-shifting transformer (PST) to regulate power flow on a monitored branch. It is structurally identical to \texttt{dctl\_ltc} but acts on the PST angle rather than on voltage: the deadband is in MW (or pu power), and the controller moves the tap to keep active power within a band around the setpoint.

\subsubsection{Logic Description}\label{model-dctl:logic-description-10}

Let \(P_{meas}\) be the active power flow on the monitored branch and \(P_{set}\) the setpoint (initialised at \(t=0\)):

\[
\text{if } |P_{meas} - P_{set}| > \delta P: \quad \text{start timer}
\]
\[
\text{if timer} \ge \tau_{delay}: \quad \text{advance/retard PST angle by } \Delta\phi
\]

The direction parameter defines whether increasing the tap increases or decreases power flow.

\subsubsection{Parameters}\label{model-dctl:parameters-10}

{\def\LTcaptype{none} % do not increment counter
\begin{small}\begin{longtable}[]{@{}lll@{}}
\toprule\noalign{}
Field & Working variable & Description \\
\midrule\noalign{}
\endhead
\bottomrule\noalign{}
\endlastfoot
\texttt{pst\_branch} & \texttt{w(1)} & PST branch name (must be type \texttt{trfo}) \\
\texttt{mon\_branch} & \texttt{w(2)} & Monitored branch for power flow \\
\texttt{direction} & \texttt{w(3)} & Sign convention for power vs.~Angle \\
\texttt{ratio\_min} & \texttt{w(4)} & Minimum PST ratio (pu) \\
\texttt{ratio\_max} & \texttt{w(5)} & Maximum PST ratio (pu) \\
\texttt{nb\_positions} & \texttt{w(6)} & Number of positions \\
\texttt{half\_deadband} & \texttt{w(7)} & Power half-deadband (pu) \\
\texttt{delay\_first} & \texttt{w(8)} & First change delay (s) \\
\texttt{delay\_next} & \texttt{w(9)} & Subsequent change delay (s) \\
\end{longtable}\end{small}
}

9 fields required. Internal variables: \texttt{w(10)} = power setpoint, \texttt{w(11)} = time of last change, \texttt{w(12)} = active delay, \texttt{w(13)} = state.

\subsubsection{Usage Example}\label{model-dctl:usage-example-10}

\begin{Verbatim}
DCTL  PST  PST_L1  PST-TR1  LINE-MON  1  -0.5  0.5  21  0.02  30.0  10.0 ;
\end{Verbatim}

\begin{center}\rule{0.5\linewidth}{0.5pt}\end{center}

\subsection{\texorpdfstring{\texttt{RT} (\texttt{dctl\_rt}): Real-Time Simulation Control}{RT (dctl\_rt): Real-Time Simulation Control}}\label{model-dctl:rt-dctl_rt-real-time-simulation-control}

\subsubsection{Description}\label{model-dctl:description-11}

Forces the RAMSES simulation to track \textbf{wall-clock time} at a configurable ratio, enabling hardware-in-the-loop (HIL) or real-time simulation testing. A ratio of 1.0 means the simulation runs in exact real time; a ratio of 2.0 allows the simulation to run up to twice as fast as real time (useful for digital twin applications).

An optional over-run tolerance parameter \texttt{allowed\_overrun} specifies how far the simulation may lag behind real time before an error is triggered.

\subsubsection{Parameters}\label{model-dctl:parameters-11}

{\def\LTcaptype{none} % do not increment counter
\begin{small}\begin{longtable}[]{@{}
  >{\raggedright\arraybackslash}p{(\linewidth - 4\tabcolsep) * \real{0.1892}}
  >{\raggedright\arraybackslash}p{(\linewidth - 4\tabcolsep) * \real{0.4595}}
  >{\raggedright\arraybackslash}p{(\linewidth - 4\tabcolsep) * \real{0.3514}}@{}}
\toprule\noalign{}
\begin{minipage}[b]{\linewidth}\raggedright
Field
\end{minipage} & \begin{minipage}[b]{\linewidth}\raggedright
Working variable
\end{minipage} & \begin{minipage}[b]{\linewidth}\raggedright
Description
\end{minipage} \\
\midrule\noalign{}
\endhead
\bottomrule\noalign{}
\endlastfoot
\texttt{ratio} & \texttt{w(1)} & Ratio to real time (1.0 = real time, \textgreater1 = faster) \\
\texttt{allowed\_overrun} & \texttt{w(2)} & Allowed over-run before error (s), optional \\
\end{longtable}\end{small}
}

1 or 2 fields. Only one \texttt{DCTL\ RT} controller may be active per simulation.

Observable: \texttt{elapsed} (wall-clock elapsed time in seconds).

\subsubsection{Usage Example}\label{model-dctl:usage-example-11}

\begin{Verbatim}
DCTL  RT  RT_CTL  1.0  0.01 ;
\end{Verbatim}

\begin{center}\rule{0.5\linewidth}{0.5pt}\end{center}

\subsection{\texorpdfstring{\texttt{dctl\_hvdc\_lim}: HVDC LCC Limiter/Controller}{dctl\_hvdc\_lim: HVDC LCC Limiter/Controller}}\label{model-dctl:dctl_hvdc_lim-hvdc-lcc-limitercontroller}

\begin{docsnote}{Note}

Callable since 3.79. Earlier releases compiled this model under no name, so a data file referring to it was rejected.

\end{docsnote}

\subsubsection{Description}\label{model-dctl:description-12}

A discrete controller that \textbf{supervises and adjusts an LCC-HVDC link} (\texttt{twop\_HVDC\_LCC}) by:
- Moving rectifier and inverter transformer tap changers to maintain firing and extinction angles within desired ranges
- Triggering reduced-current or blocked operation when AC voltages are too low for normal commutation
- Setting the \texttt{alphamin}, \texttt{gammamin}, and \texttt{Idred} parameters of the HVDC model

It interacts directly with the HVDC two-port model's internal parameters.

\subsubsection{Logic Description}\label{model-dctl:logic-description-11}

\textbf{Rectifier tap changer}: if \(\alpha\) is outside \([\alpha_{min,des}, \alpha_{max,des}]\) for longer than \texttt{delay\_rect}, move rectifier transformer tap.

\textbf{Inverter tap changer}: if \(\gamma\) is outside \([\gamma_{min,des}, \gamma_{max,des}]\) for longer than \texttt{delay\_inv}, move inverter transformer tap.

\textbf{Reduced-current control}: when \(V_{dc}\) falls below \(V_{min,des}\) or \(V_{max,des}\) the controller switches the inverter to reduced-current mode (\texttt{inv\_ctl\ =\ -1}) and sets \(I_{red}\).

\textbf{Blocked operation}: if conditions are unrecoverable (very low AC voltage), the HVDC is blocked (\texttt{rec\_ctl\ =\ 3}, \texttt{inv\_ctl\ =\ 3}).

\subsubsection{Parameters}\label{model-dctl:parameters-12}

{\def\LTcaptype{none} % do not increment counter
\begin{small}\begin{longtable}[]{@{}
  >{\raggedright\arraybackslash}p{(\linewidth - 4\tabcolsep) * \real{0.1892}}
  >{\raggedright\arraybackslash}p{(\linewidth - 4\tabcolsep) * \real{0.4595}}
  >{\raggedright\arraybackslash}p{(\linewidth - 4\tabcolsep) * \real{0.3514}}@{}}
\toprule\noalign{}
\begin{minipage}[b]{\linewidth}\raggedright
Field
\end{minipage} & \begin{minipage}[b]{\linewidth}\raggedright
Working variable
\end{minipage} & \begin{minipage}[b]{\linewidth}\raggedright
Description
\end{minipage} \\
\midrule\noalign{}
\endhead
\bottomrule\noalign{}
\endlastfoot
\texttt{hvdc\_name} & \texttt{w(1)} & Target HVDC two-port device name \\
\texttt{nr\_min} & \texttt{w(2)} & Minimum rectifier transformer ratio \\
\texttt{nr\_max} & \texttt{w(3)} & Maximum rectifier transformer ratio \\
\texttt{nr\_step} & \texttt{w(4)} & Rectifier tap step size \\
\texttt{ni\_min} & \texttt{w(5)} & Minimum inverter transformer ratio \\
\texttt{ni\_max} & \texttt{w(6)} & Maximum inverter transformer ratio \\
\texttt{ni\_step} & \texttt{w(7)} & Inverter tap step size \\
\texttt{alpha\_min\_des} & \texttt{w(8)} & Desired minimum firing angle (degrees) \\
\texttt{alpha\_max\_des} & \texttt{w(9)} & Desired maximum firing angle (degrees) \\
\texttt{gamma\_min\_des} & \texttt{w(10)} & Desired minimum extinction angle (degrees) \\
\texttt{gamma\_max\_des} & \texttt{w(11)} & Desired maximum extinction angle (degrees) \\
\texttt{V\_min\_des} & \texttt{w(12)} & Minimum desired DC voltage (kV) \\
\texttt{V\_max\_des} & \texttt{w(13)} & Maximum desired DC voltage (kV) \\
\texttt{delay\_rect} & \texttt{w(14)} & Delay between rectifier tap changes (s) \\
\texttt{delay\_inv} & \texttt{w(15)} & Delay between inverter tap changes (s) \\
\texttt{alphamin} & \texttt{w(16)} & Minimum firing angle limit (degrees) \\
\texttt{Idred} & \texttt{w(17)} & Reduced DC current (kA) \\
\texttt{gammamin} & \texttt{w(18)} & Minimum extinction angle (degrees) \\
\end{longtable}\end{small}
}

18 fields required. Internal variables: \texttt{w(19)}--\texttt{w(22)} (operation modes and tap-change timers).

\subsubsection{Usage Example}\label{model-dctl:usage-example-12}

\begin{Verbatim}
DCTL  HVDC_LIM  HVDCLIM1  HVDClink
      0.85  1.15  0.00625
      0.85  1.15  0.00625
      15.0  18.0  15.0  18.0
      400.0  600.0
      10.0  10.0
      5.0  0.0  15.0 ;
\end{Verbatim}

\begin{center}\rule{0.5\linewidth}{0.5pt}\end{center}

\section{Monitoring}\label{model-dctl:monitoring}

\subsection{\texorpdfstring{SIM\_MINMAXSPEED (\texttt{dctl\_sim\_minmaxspeed}): Speed Monitoring (Min/Max)}{SIM\_MINMAXSPEED (dctl\_sim\_minmaxspeed): Speed Monitoring (Min/Max)}}\label{model-dctl:sim_minmaxspeed-dctl_sim_minmaxspeed-speed-monitoring-minmax}

\subsubsection{Description}\label{model-dctl:description-13}

A \textbf{simulation monitoring} controller that checks whether any synchronous machine speed violates upper or lower bounds. If a violation persists beyond a deadtime, the controller logs a warning or stops the simulation, depending on the \texttt{stop\_sim} flag. It is used to automatically terminate simulations that have become numerically unstable or that have reached an unacceptable operating state.

\subsubsection{Logic Description}\label{model-dctl:logic-description-12}

At each evaluation, the controller scans all synchronous machines and checks:

\[
\text{if } \omega > \omega_{max} \text{ or } \omega < \omega_{min}: \quad \text{record violation time}
\]
\[
\text{if violation duration} > T_{deadtime}: \quad \text{warn or stop simulation}
\]

\subsubsection{Parameters}\label{model-dctl:parameters-13}

{\def\LTcaptype{none} % do not increment counter
\begin{small}\begin{longtable}[]{@{}
  >{\raggedright\arraybackslash}p{(\linewidth - 4\tabcolsep) * \real{0.1892}}
  >{\raggedright\arraybackslash}p{(\linewidth - 4\tabcolsep) * \real{0.4595}}
  >{\raggedright\arraybackslash}p{(\linewidth - 4\tabcolsep) * \real{0.3514}}@{}}
\toprule\noalign{}
\begin{minipage}[b]{\linewidth}\raggedright
Field
\end{minipage} & \begin{minipage}[b]{\linewidth}\raggedright
Working variable
\end{minipage} & \begin{minipage}[b]{\linewidth}\raggedright
Description
\end{minipage} \\
\midrule\noalign{}
\endhead
\bottomrule\noalign{}
\endlastfoot
\texttt{omega\_max} & \texttt{w(1)} & Upper speed limit (pu) \\
\texttt{omega\_min} & \texttt{w(2)} & Lower speed limit (pu) \\
\texttt{deadtime} & \texttt{w(3)} & Deadtime before action (s) \\
\texttt{stop\_sim} & \texttt{w(4)} & 1 = stop simulation on violation, 0 = warn only \\
\end{longtable}\end{small}
}

4 fields required. Internal: \texttt{w(5)} = last violation time.

Parameter names exported: \texttt{SPEEDMIN}, \texttt{SPEEDMAX}, \texttt{DEADTIME}.

\subsubsection{Usage Example}\label{model-dctl:usage-example-13}

\begin{Verbatim}
DCTL  SIM_MINMAXSPEED  SPEEDMON  1.2  0.5  0.5  1 ;
\end{Verbatim}

\begin{center}\rule{0.5\linewidth}{0.5pt}\end{center}

\subsection{\texorpdfstring{SIM\_MINMAXVOLT (\texttt{dctl\_sim\_minmaxvolt}): Voltage Monitoring (Min/Max)}{SIM\_MINMAXVOLT (dctl\_sim\_minmaxvolt): Voltage Monitoring (Min/Max)}}\label{model-dctl:sim_minmaxvolt-dctl_sim_minmaxvolt-voltage-monitoring-minmax}

\subsubsection{Description}\label{model-dctl:description-14}

Analogous to \texttt{dctl\_sim\_minmaxspeed} but for \textbf{bus voltages}. The controller scans all buses and checks whether any voltage violates the specified bounds. If the violation persists for longer than the deadtime, a warning is logged or the simulation is terminated.

\subsubsection{Logic Description}\label{model-dctl:logic-description-13}

\[
\text{if } V > V_{max} \text{ or } V < V_{min}: \quad \text{record violation time}
\]
\[
\text{if violation duration} > T_{deadtime}: \quad \text{warn or stop simulation}
\]

\subsubsection{Parameters}\label{model-dctl:parameters-14}

{\def\LTcaptype{none} % do not increment counter
\begin{small}\begin{longtable}[]{@{}lll@{}}
\toprule\noalign{}
Field & Working variable & Description \\
\midrule\noalign{}
\endhead
\bottomrule\noalign{}
\endlastfoot
\texttt{V\_max} & \texttt{w(1)} & Upper voltage limit (pu) \\
\texttt{V\_min} & \texttt{w(2)} & Lower voltage limit (pu) \\
\texttt{deadtime} & \texttt{w(3)} & Deadtime before action (s) \\
\texttt{stop\_sim} & \texttt{w(4)} & 1 = stop on violation, 0 = warn only \\
\end{longtable}\end{small}
}

4 fields required. Internal: \texttt{w(5)} = last violation time.

Parameter names exported: \texttt{VMIN}, \texttt{VMAX}, \texttt{DEADTIME}.

\subsubsection{Usage Example}\label{model-dctl:usage-example-14}

\begin{Verbatim}
DCTL  SIM_MINMAXVOLT  VOLTMON  1.15  0.80  1.0  0 ;
\end{Verbatim}

\begin{center}\rule{0.5\linewidth}{0.5pt}\end{center}

\subsection{\texorpdfstring{VOLT\_VAR (\texttt{dctl\_voltage\_variability}): Voltage Variability Monitor}{VOLT\_VAR (dctl\_voltage\_variability): Voltage Variability Monitor}}\label{model-dctl:volt_var-dctl_voltage_variability-voltage-variability-monitor}

The data-file model name is \texttt{VOLT\_VAR}.

\subsubsection{Description}\label{model-dctl:description-15}

Computes a \textbf{running Exponential Moving Average (EMA)} of voltage variability (variance) across a list of monitored buses. It is used to assess power quality metrics such as voltage flicker or ramping caused by variable renewable energy sources.

The EMA uses two decay constants (\(\lambda_1\), \(\lambda_2\)) derived from the averaging time window, enabling a computationally efficient recursive estimation:

\[
\text{EMA}(t) = \lambda_1 \cdot \text{EMA}(t-1) + (1 - \lambda_1) \cdot V(t) + (\lambda_1 - \lambda_2)(V(t) - V(t-1))
\]

The controller accumulates per-bus variances and a system-wide sum.

\subsubsection{Parameters}\label{model-dctl:parameters-15}

{\def\LTcaptype{none} % do not increment counter
\begin{small}\begin{longtable}[]{@{}
  >{\raggedright\arraybackslash}p{(\linewidth - 2\tabcolsep) * \real{0.3500}}
  >{\raggedright\arraybackslash}p{(\linewidth - 2\tabcolsep) * \real{0.6500}}@{}}
\toprule\noalign{}
\begin{minipage}[b]{\linewidth}\raggedright
Field
\end{minipage} & \begin{minipage}[b]{\linewidth}\raggedright
Description
\end{minipage} \\
\midrule\noalign{}
\endhead
\bottomrule\noalign{}
\endlastfoot
\texttt{aver\_time\_window} & Averaging time window (s) for EMA computation \\
\texttt{nb\_list} & Number of buses to monitor (bus names follow in subsequent RAMSES data records) \\
\end{longtable}\end{small}
}

2 fields required. Only one \texttt{VOLT\_VAR} controller may exist per simulation. Parameters are stored in module-level variables (\texttt{volt\_var\_mod}) rather than \texttt{w(*)}.

\subsubsection{Usage Example}\label{model-dctl:usage-example-15}

\begin{Verbatim}
DCTL  VOLT_VAR  VVAR  60.0  5 ;
\end{Verbatim}
